\subsection{正三次元側における商体への退化の排除}
\label{subsec:h3-positive3-rational-collapse}
\begin{lemma}\label{lem:h3-positive3-rational-collapse}
重複度$(3,1)$のH3包絡で、共通平面の式が$q_0(W)=0$、すなわち
系\ref{cor:h3-order2-plane-quadric}の$\gamma=0$であるとする。
$T=\langle W\rangle$をそのAS2平面、$z$を負生成元、$v$を残りの正生成元とする。
$T$の斉次商体に$v$と$z^2$がそれぞれ次数1、2の元として同時に入ることはない。
従って、二次関係から$v\ell=q(W)$または$\ell v=q(W)$と
$z^2$を$v,W$で表す二次式が得られる場合は排除される。
単冪部分を零と仮定しない。$\gamma\ne0$にはこの補題を適用しない。
\end{lemma}
\begin{proof}
$D=Q(T)$とし、仮に$v,h=z^2\in D$とする。
$zt=f(t)z$、$\det f=\operatorname{hdet}_T f=1$、
$\mu_T=\mu_B|_W\circ f^2$が既に証明されている。
$h\ne0$は$ht=f^2(t)h$を満たす。

$T$が可換なら、内作用$\operatorname{Ad}_h$は恒等なので$f^2=1$。
従って$f$は有限位数であり、$D$は固定体$F=D^{\langle f\rangle}$上有限次元、
$\operatorname{trdeg}_kF=2$である。
$B\subset D+Dz$も$F$上有限次元の代数に入り、$F$はこの代数の中心に含まれる。
有限個の生成元係数と構造定数を含む有限生成可換部分環を$F$内に取ると、
長さ$n$の語の係数次数は$O(n)$に抑えられる。
従って$\operatorname{GKdim}B\le2$となり矛盾する。

$T$が量子平面$xy=qyx$、$q\ne1$なら、$\mu_B|_W$の固有値は1だけなので
その作用は恒等である。$q=-1$の反対角自己同型は
$\operatorname{hdet}f=-\det f$となるので今の条件では存在しない。
従って$f=\operatorname{diag}(a,a^{-1})$、$a^2=q$である。
$q$が有限位数なら$f$も有限位数で、$D$は二変数中心体上有限次元である。
その中心の$f$固定体を$F$として、前段と同じ有限構造定数の議論によりGK次元2以下の矛盾を得る。
$q$が無限位数なら、斉次中心元は$k$だけである。
実際、$u=xy^{-1}$により斉次元を$r(u)y^m$と書くと、
$y$との可換性は$r(q^{-1}u)=r(u)$を要求する。
有理関数の有限な零極が無限軌道を持ち得ないため$r$は定数、
さらに$x$との可換性から$q^m=1$、従って$m=0$となる。
一方、$f^2=\mu_T$は$(xy)^{-1}$による内作用なので、
$hxy$は中心に入り、非零次数4である。
これは中心が次数零の$k$であることに反する。

残るJordan平面を$yx=xy+y^2$と書く。
自己同型は$f(y)=a y,f(x)=a x+b y$であり、
$\operatorname{hdet}f=a^2=1$から$f^2(y)=y$である。
従って$hy=yh$。
$u=xy^{-1}$とすると$yu=(u+1)y$で、次数2の$h$は
$h=r(u)y^2$、$r\in k(u)^\times$と書ける。
可換条件は$r(u+1)=r(u)$となる。
有限の極は平行移動の無限軌道を作るため存在せず、
周期多項式も標数0では定数だけなので$h=c y^2$、$c\ne0$である。
すると元の整域の商体内で
\[
 (zy^{-1})^2=ac\in k^\times.
\]
代数閉体上ではこの多項式が一次式へ分解するため、整域性から$zy^{-1}\in k$。
従って$z\in ky$となり、一次元負固有空間と正固有空間の独立性に反する。
これで全AS2平面を扱った。
最後の主張では$\ell\ne0$の逆元で$v$が商体に入り、
外側の平方関係から$z^2$も同じ商体に入るため、証明した矛盾を適用できる。
\end{proof}

この補題は、正三次元側でcross射影が退化した場合を、
既存の負三次元側と同様に排除するために使える。
ただし$\gamma=0$、外側平方関係、cross射影の次元等を先に確認する必要があり、
全normalizing chartのH3排除や全$(3,1)$被覆はまだ従わない。
